\summarychapter{22}{Sleeping Men and the Law of Necessity}{The Decline of the Papacy}

\begin{inbrief}
The fourteenth century opened with a triumphant Jubilee and within three years saw the greatest
insult ever offered a pope. The chapter follows the fall of papal sovereignty: the defeat of Boniface
VIII by the new national monarchies, seventy-two years of popes at Avignon, the Great Schism with two
and then three popes, and the councils that ended the schism and claimed authority above the pope.
\question{What made rulers and peoples turn against the popes' claim to worldly power?}
\end{inbrief}

On February 22, 1300, Boniface VIII (1294--1303) proclaimed the first Holy Year, promising ``a full
and copious pardon'' to all who came to the churches of St.~Peter and of St.~Paul; Paul VI echoed him in
announcing the Holy Year of 1975. Pilgrims poured in, and at St.~Paul's two priests were kept busy day
and night raking together the offerings. Boniface loved pomp. He appeared in imperial robes crying,
``I am Caesar, I am emperor,'' and wore a crown heavy with jewels. With the example of Innocent III
before him he expected to command kings as his predecessors had.

\ssection{Sleep and Change in the Kingdom}
In Jesus' parable the seed grows night and day while the sower sleeps and does not know how. Great
changes often happen unnoticed, and so it was in the fourteenth century: the papacy carried on as usual while new ideas and social forces altered Christianity's face. The years 1300 to
1500 are called the decline of the Middle Ages, since Christendom, the idea that had held a thousand years together, now came under attack. Christendom had two pillars, a Christian empire and a Catholic Church. The empire had already faded as the papacy rose; now papal prestige fell as well.
People were getting used to thinking of themselves as English or French and of their state as working
without papal guidance---in modern terms, of a secular realm alongside a sacred one.

Europe was leaving feudalism behind. Cash counted for more than land, and rulers needed wider powers
to tax. Edward I of England and Philip the Fair of France, at war over English lands in France, both
decided to tax the clergy. In 1296 Boniface answered with \emph{Clericis laicos}, threatening to
excommunicate any ruler who taxed clergy and any cleric who paid without papal consent. Edward
replied by withdrawing legal protection from clergy who refused; Philip forbade the export of gold,
silver and jewels, cutting off the papal income from France. Boniface retreated and explained that he had never intended to forbid contributions for defense in dire need. Since the kings decided what that
meant, they had won.

\ssection{A New Voice in Christendom}
Encouraged by the Jubilee, Boniface took the reverence shown him as reaching to civil matters too and
added a second circlet to his crown as a sign of temporal sovereignty. Philip represented a new voice:
Christ had left the church no power in worldly affairs. In 1301 the king had a French bishop jailed for treason. Boniface demanded his release and withdrew his concession on taxes. Philip summoned nobles,
clergy and townsmen in 1302---the first Estates-General---and won their unanimous support. ``My
master's sword is made of steel; the pope's is made of words,'' said one of his ministers. Boniface
then issued \emph{Unam sanctam}, the most extreme claim of papal power in history: ``It is altogether
necessary for every human being to be subject to the Roman pontiff.''

Philip prepared to have him deposed on the ground that his election was illegal, adding charges of
heresy, simony and immorality, and sent the lawyer William of Nogaret, a master of the trumped-up
charge, to bring him to France for trial. Nogaret and his troops broke in on the pope, now eighty-six,
in his bedroom at Anagni, his birthplace, abused him and held him for several days until the
townspeople rescued him. Humiliated, Boniface died within weeks. ``He crept in like a fox, reigned
like a lion, and died like a dog.'' Anagni showed that Europeans no longer accepted papal
interference in what they considered political matters, even though the outrage was widely resented.
Boniface was not loved; Dante put him in hell. Yet few could imagine Christianity without the pope.
People were beginning to distinguish secular from religious authority and to grant each its rights.
As Canossa had marked the rise of papal power, Anagni marked its descent. In 1305 the cardinals chose the French archbishop of Bordeaux, who became Clement~V. He never set foot in Rome.

\ssection{The Captivity of the Holy Father}
So began the seventy-two years called the Babylonian Captivity of the papacy. After Clement six more
popes, all French, lived at Avignon on the Rhone, just across from Philip's lands, which grew into a
city of eighty thousand with a huge bureaucracy and a sumptuous palace. Rome stood for the succession
from Peter and for western universality; Avignon looked like the tool of one power-hungry nation.
Germany resented it bitterly. In 1324 the emperor Louis the Bavarian appealed against John XXII to a
general council. Marsilius of Padua, who had fled the University of Paris, and John of Jandun
presented him in 1326 with the \emph{Defender of the Peace}. It taught that the church is the
community of all believers and the priesthood is not superior to the laity. Popes, bishops and priests
have no special function from Christ but are agents of the community, which is represented by the
general council. The pope becomes the council's executive officer. This theory is conciliarism.

Most anger at Avignon concerned money. With revenue from the Papal States fallen, the popes multiplied
fees and taxes: the annat, the first year's income of a newly appointed bishop, collected more often
by transferring bishops from see to see; the reservation, by which the pope left an office vacant and
kept its income; and above all indulgences, granted for the slightest reasons with ever larger
claims, while payment was demanded under threat of excommunication. By 1360 it was clear that Avignon
could not last. In 1377 the aged Gregory XI returned to Rome and died within a year. The cardinals,
still mostly French, gave way to a Roman mob and chose an Italian, Urban VI, crowned on Easter Sunday,
April 18, with all of them present. His dictatorial ways brought second thoughts. In August they
declared the election forced and invalid. Urban created an almost new college of cardinals; the
French cardinals elected one of their own, Clement VII, who went back to Avignon.

\ssection{The Great Papal Schism}
For thirty-nine years there were two lines of popes, each with its own cardinals, each claiming to be
the true Vicar of Christ and excommunicating the other's followers. There had been schisms before, but
now the Church of Rome itself was both promoting and rejecting its own pontiff. The only witnesses,
the cardinals, flatly contradicted themselves, so that universities, kings and bishops had to decide.
France and Scotland followed Clement; Italy, the empire and England followed Urban; every country had
minorities. There were riots, burnings and even crusades. In 1395 professors at the University of
Paris proposed a general council. But canon law said that only a pope could call a council or ratify
its decisions---and which pope? In effect the law forbade Christendom's reunion. Did necessity rank above the law?

\ssection{A New Pope Elected}
In 1409 most cardinals of both sides decided that it was. Meeting at Pisa, they deposed both popes
and elected Alexander V. Neither deposed pope gave way, so there were now three, one of whom preached
a crusade against another and sold indulgences to pay for it. The spectacle forced action. In 1414
the Holy Roman Emperor gathered the most impressive council of the age at Constance, with even Greek
Orthodox representatives. For the first time voting was by ``nations''---German, Italian, French,
English and later Spanish, each with one vote---and laymen took part, a sign that the church was
reluctantly recognizing the new national alignment of power. (A chart in the chapter is headed ``The
Decline of the Middle Ages.'') In 1417 the council persuaded one pope to resign, deposed the other
two and elected Martin V. Benedict XIII at Avignon would not give up his claim, but the schism was over.

Necessity had triumphed and was at once denied. Martin, who reigned only by the council's act,
repudiated all its acts except his own election, for the council raised the question whether a
council that makes a pope is greater than the pope. The conciliar movement wanted a limited monarchy;
Constance decreed that general councils are superior to popes and must meet regularly. The popes
called this heresy, and since later councils failed to bring reform, they had discredited the
movement by 1450. They then busied themselves with Italian politics and patronage of the arts, unsure
whether they were successors of Peter or of Caesar. Corruption reached its depth under Roderigo
Borgia, Alexander VI (1492--1503), grossly immoral and bent on enriching his children. Constance
could be denied but not forgotten. People had begun to think of national churches and of a church
governed by representative bodies. The ground for the Protestant Reformation was prepared.

\begin{keyfacts}
\keypeople{Boniface VIII; Philip the Fair; Edward I; William of Nogaret; Dante; Clement V; Louis the
Bavarian; John XXII; Marsilius of Padua; John of Jandun; Gregory XI; Urban VI; Clement VII; Alexander
V; Martin V; Benedict XIII; Alexander VI (Roderigo Borgia).}
\keydates{%
  \kd{1296}{\emph{Clericis laicos}: Boniface VIII forbids taxing the clergy}
  \kd{1300}{Boniface VIII proclaims the first Jubilee (Holy Year)}
  \kd{1302}{Boniface VIII's \emph{Unam sanctam}; Philip the Fair calls the first Estates-General}
  \kd{1303}{Boniface VIII seized at Anagni; he dies within weeks}
  \kd{1305}{Clement V elected: the Avignon papacy (``Babylonian Captivity'') begins}
  \kd{1326}{Marsilius of Padua's \emph{Defender of the Peace}}
  \kd{1377}{Gregory XI returns the papacy to Rome}
  \kd{1378}{Two popes elected: the Great Schism begins}
  \kd{1409}{Council of Pisa elects a third pope}
  \kd{1414--1417}{Council of Constance ends the schism; Martin V elected}
  \kd{1492--1503}{Pontificate of Alexander VI (Roderigo Borgia)}}
\keyterms{Jubilee; plenary indulgence; \emph{Clericis laicos}; \emph{Unam sanctam}; Estates-General;
Babylonian Captivity; Avignon; annat; reservation; Great Schism; conciliarism; council of ``nations.''}
\end{keyfacts}
